\section{Introduction}

\subsection{Motivation and Background}

The telecommunications industry is undergoing a fundamental shift from
purpose-built, vertically integrated network appliances towards software
running on commodity computing infrastructure. This transformation,
formalised by ETSI as \textbf{Network Functions Virtualisation (NFV)},
decouples network functions (firewalls, routers, load balancers, packet
gateways, etc.) from the proprietary hardware on which they traditionally
ran \cite{etsi002}. In parallel, \textbf{Software-Defined Networking (SDN)}
separates the network control plane from the data plane, enabling
centralised, programmable network management.

The first generation of NFV deployments virtualised network functions as
Virtual Machines (VMs), managed by OpenStack as the Virtual Infrastructure
Manager (VIM). More recently the industry has embraced \textbf{cloud-native}
principles: network functions are packaged as containers and orchestrated by
Kubernetes, giving rise to \textbf{Cloud-Native Network Functions (CNFs)}
\cite{etsi022,lfn}. CNFs boot in seconds, consume far fewer resources, and
inherit Kubernetes' self-healing and horizontal scaling capabilities.

\subsection{Problem Statement}

Production NFV platforms are large, expensive, and operationally complex,
which makes them impractical for learning and experimentation. This project
addresses the question: \emph{can the essential structure and behaviour of
the ETSI NFV framework be reproduced faithfully on a single personal
computer?} The challenge is to build a self-contained testbed that
(i)~maps cleanly onto the ETSI reference architecture, (ii)~exhibits the
telco-grade separation of management, control, and user planes, and
(iii)~fits within a hard budget of 16~GB of host RAM.

\subsection{Objectives}

\begin{enumerate}
    \item Provision a multi-node Kubernetes cluster entirely through
          Infrastructure-as-Code (Vagrant~+~Ansible).
    \item Deploy a CNF with \textbf{multi-NIC} networking that separates
          management, cluster, and data-plane traffic.
    \item Reproduce the telco three-plane network model using a secondary
          CNI (Multus~+~macvlan~+~Whereabouts).
    \item Provide full observability with Prometheus and Grafana.
    \item Demonstrate horizontal auto-scaling (HPA) and measure real
          data-plane throughput.
    \item Establish an explicit mapping between each ETSI NFV building block
          and its Kubernetes equivalent.
\end{enumerate}

\subsection{Report Organisation}

Section~2 reviews the theoretical background (NFV, SDN, OpenStack, and
Kubernetes as a cloud-native VIM). Section~3 discusses the migration from
VM-based VNFs to CNFs and the multi-NIC networking model. Section~4 presents
the system architecture. Section~5 details the implementation. Section~6
reports the testing and evaluation across five test cases. Section~7
concludes and outlines future work.
